\documentclass[conference]{IEEEtran}

\IEEEoverridecommandlockouts

\usepackage[T1]{fontenc}
\usepackage[utf8]{inputenc}
\usepackage{cite}
\usepackage{amsmath,amssymb,amsfonts}
\usepackage{algorithmic}
\usepackage{graphicx}
\usepackage{textcomp}
\usepackage{xcolor}
\usepackage{booktabs}
\usepackage{multirow}
\usepackage{url}

\graphicspath{{figures/}}

\newcommand{\todo}[1]{\textcolor{red}{[TODO: #1]}}
\newcommand{\placeholderfigure}[2]{%
  \begin{figure}[!t]
    \centering
    \IfFileExists{figures/#1}{%
      \includegraphics[width=\linewidth]{#1}%
    }{%
      \fbox{\begin{minipage}[c][0.25\textheight][c]{0.92\linewidth}
      \centering #2\\[0.5em]\path{figures/#1}
      \end{minipage}}%
    }
    \caption{#2}
    \label{fig:#1}
  \end{figure}
}

\begin{document}

\title{Physics-Constrained Phase Retrieval with AutoPhaseNN and UMamba-Based Networks}

\author{
\IEEEauthorblockN{Author Name}
\IEEEauthorblockA{
Affiliation\\
City, Country\\
email@example.com}
}

\maketitle

\begin{abstract}
This paper investigates neural phase retrieval for coherent diffraction imaging
through a controlled reproduction and model-replacement study. We first
reproduce an AutoPhaseNN-style pipeline, then replace the original convolutional
backbone with a UMamba-based model while keeping the data flow, Fourier-domain
post-processing, training objective, checkpointing, and evaluation protocol
aligned. The experiments reveal a practical failure mode: far-field fitting
metrics can improve while the reconstructed real-space amplitude, phase, and
support remain physically implausible. We use this gap to analyze which parts of
the AutoPhaseNN pipeline provide effective spatial priors and what must be
controlled before attributing performance differences to the backbone itself.
\end{abstract}

\begin{IEEEkeywords}
phase retrieval, coherent diffraction imaging, AutoPhaseNN, UMamba, physics-informed deep learning
\end{IEEEkeywords}

\section{Introduction}

Phase retrieval aims to recover a complex-valued object from diffraction
measurements whose phase information is not directly observed. In coherent
diffraction imaging, this inverse problem is highly non-linear and admits
ambiguities unless additional physical priors are imposed.

Deep neural networks offer a data-driven alternative to iterative phase
retrieval algorithms. AutoPhaseNN demonstrates that a convolutional neural
network can be trained to map diffraction amplitudes to object-domain amplitude
and phase, followed by a differentiable Fourier-domain reconstruction step.
In this work, AutoPhaseNN is used first as a reproduction target rather than
only as a baseline. This choice is important: before testing a new architecture,
the data preprocessing, Fourier projection, support masking, checkpoint loading,
and evaluation metrics must match the reference pipeline closely enough that
model differences are not confused with implementation differences.

The current experimental evidence suggests a central problem: optimizing
diffraction-domain L1 loss alone does not guarantee a physically meaningful
real-space solution. This motivates a careful comparison between Fourier-domain
metrics and real-space metrics such as amplitude error, phase error, and support
agreement.

The main question is therefore not simply whether a UMamba-based network can
reduce far-field loss. The more useful question is whether it can do so while
recovering an object-domain representation that is consistent with the physical
support and phase structure expected by the phase retrieval problem. This paper
reports the current reproduction results, the UMamba replacement experiments,
and the diagnostic evidence behind this distinction.

The contributions of this draft are:
\begin{itemize}
  \item We establish a reproducible AutoPhaseNN-style PyTorch training and
  evaluation pipeline.
  \item We compare the original AutoPhaseNN backbone and UMamba-based variants
  under aligned data loading, post-processing, checkpointing, and metrics.
  \item We show how far-field losses can be inconsistent with real-space
  reconstruction quality, especially when support structure is weakly enforced.
  \item We organize the evidence for spatial support priors as a key factor in
  stable neural phase retrieval.
\end{itemize}

\section{Related Work}

\subsection{Classical Phase Retrieval}
Classical phase retrieval methods usually combine Fourier-domain consistency
with object-domain constraints such as finite support, positivity, and
oversampling assumptions~\cite{gerchberg1972practical,fienup1982phase}.
\todo{Add references to HIO and oversampling theory.}

\subsection{AutoPhaseNN}
AutoPhaseNN trains a neural network to predict object-domain amplitude and phase
from diffraction measurements. The predicted amplitude and phase are converted
into a complex object, masked by a support constraint, and transformed back to
the far field with FFT~\cite{autophasenn}. The training loss is then computed
in the diffraction domain.

\subsection{State Space Models and UMamba}
Mamba and UMamba introduce state-space sequence modeling into vision and medical
image segmentation architectures~\cite{gu2023mamba,umamba}. UMamba-style blocks
may provide a larger effective receptive field than local convolution alone, but
the architecture does not necessarily encode the same hard spatial support prior
as the original AutoPhaseNN decoder.

\section{Method}

\subsection{Problem Formulation}
Let $x \in \mathbb{C}^{N \times N \times N}$ denote the object-domain complex
volume. The measured diffraction modulus is
\begin{equation}
  y = \left|\mathcal{F}\{x\}\right|,
\end{equation}
where $\mathcal{F}$ denotes the 3D Fourier transform. The network receives $y$
and predicts object-domain amplitude $\hat{a}$ and phase $\hat{\phi}$.

\subsection{Network Output and Post-processing}
The predicted complex object is constructed as
\begin{equation}
  \hat{x} = s(\hat{a}) \cdot \hat{a}\exp(i\hat{\phi}),
\end{equation}
where $s(\hat{a})$ is the support mask. In the AutoPhaseNN-style pipeline, a
hard support threshold is applied:
\begin{equation}
  s(\hat{a}) =
  \begin{cases}
  1, & \hat{a} \geq T,\\
  0, & \hat{a} < T.
  \end{cases}
\end{equation}
The predicted diffraction modulus is
\begin{equation}
  \hat{y} = \left|\mathcal{F}\{\operatorname{ifftshift}(\hat{x})\}\right|.
\end{equation}
This layer is central to the comparison because it couples three quantities:
the predicted amplitude determines the support mask, the support mask gates the
complex object, and the gated object determines the far-field loss. A low
far-field error therefore does not directly imply that $\hat{a}$, $\hat{\phi}$,
or $s(\hat{a})$ are individually correct.

\subsection{Training Objective}
The main training objective follows the paper-style far-field modulus L1 loss:
\begin{equation}
  \mathcal{L}_{FT} = \frac{1}{|\Omega|}\sum_{v \in \Omega}
  \left| \hat{y}_v - y_v \right|.
\end{equation}
Additional real-space metrics are reported for diagnosis, but the default
AutoPhaseNN reproduction uses only the diffraction-domain loss.

\subsection{Controlled Backbone Replacement}
The UMamba experiment is designed as a controlled replacement rather than a new
pipeline. The intended controlled variables are:
\begin{itemize}
  \item identical train/validation sample selection;
  \item identical preprocessing of diffraction, amplitude, and phase tensors;
  \item identical physical post-processing from amplitude and phase to far-field
  modulus;
  \item identical far-field training loss unless an ablation explicitly changes
  it;
  \item identical evaluation metrics and visualization protocol.
\end{itemize}
Under this setup, any persistent difference in reconstruction quality is more
likely to come from the model architecture or from the strength of the encoded
spatial prior.

\subsection{Model Variants}
We compare the following models:
\begin{itemize}
  \item AutoPhaseNN reproduction with the original convolutional architecture.
  \item UMamba replacement with the same data and post-processing pipeline.
  \item UMamba variants with modified support thresholding or explicit spatial
  priors.
\end{itemize}

\section{Experiments}

\subsection{Dataset}
The experiments use simulated 3D diffraction data and corresponding real-space
complex volumes. The training split contains \todo{number} samples and the
validation/test split contains \todo{number} samples. Each sample has spatial
size $64^3$.

\subsection{Training Setup}
Unless otherwise specified, models are trained from scratch with Adam and an
initial learning rate of \todo{learning rate}. The default loss is far-field L1.
For small-sample debugging, we also use an overfitting protocol in which a fixed
subset of samples is used for both training and validation.

The small-sample protocol is used to answer a narrower question than general
validation accuracy: can the model fit a known fixed subset and produce
physically plausible object-domain amplitude and phase for those same samples?
This prevents generalization error from being confused with an optimization or
architecture-prior problem.

\subsection{Experiment Design}
The paper focuses on a small number of experiments that directly support the
main claim, rather than reporting every debugging run. The core comparison is:
\begin{itemize}
  \item AutoPhaseNN reproduction from scratch with the paper-style far-field
  objective.
  \item UMamba backbone replacement with the same data interface,
  post-processing, objective, and evaluation protocol.
  \item One representative diagnostic ablation, selected only if it helps
  explain why Fourier-domain fitting and real-space correctness diverge.
\end{itemize}
Learning-rate sweeps, checkpoint debugging, gradient probes, and implementation
sanity checks are treated as internal controls and are not reported as main
experiments unless they change the scientific conclusion.

\subsection{Evaluation Metrics}
We report fixed metric groups:
\begin{itemize}
  \item FT metrics: L1, MSE, RMSE, and relative L1 of diffraction modulus.
  \item Amplitude metrics: full-volume L1, MSE, RMSE, and relative L1.
  \item Phase metrics: wrapped phase error on the true support.
  \item Support metrics: L1, IoU, Dice score, and predicted support fraction.
\end{itemize}

\subsection{Implementation Details}
The AutoPhaseNN and UMamba pipelines share the same dataset interface,
post-processing logic, checkpoint format, and evaluation scripts wherever
possible. This controls for non-model factors in the comparison.

\section{Results}

\subsection{Visual Summary}
The first result to present is visual rather than only numerical: compare the
reconstructed diffraction modulus, amplitude, phase, and support for the same
sample IDs across model variants. This directly shows whether a lower
far-field metric corresponds to a plausible object-domain reconstruction.

\placeholderfigure{fig_autophasenn_reproduction.png}{AutoPhaseNN reproduction visualization.}

\placeholderfigure{fig_umamba_original.png}{UMamba original hard-threshold visualization.}

\placeholderfigure{fig_umamba_soft_threshold.png}{UMamba soft-threshold visualization.}

\placeholderfigure{fig_training_curves.png}{Training and validation loss curves from TensorBoard event files.}

\begin{table}[!t]
  \caption{Evaluation Metrics of Model Variants}
  \label{tab:main_results}
  \centering
  \begin{tabular}{lcccc}
    \toprule
    Model & FT L1 $\downarrow$ & Amp L1 $\downarrow$ & Phase L1 $\downarrow$ & Support IoU $\uparrow$ \\
    \midrule
    AutoPhaseNN & \todo{} & \todo{} & \todo{} & \todo{} \\
    UMamba & \todo{} & \todo{} & \todo{} & \todo{} \\
    Diagnostic ablation & \todo{} & \todo{} & \todo{} & \todo{} \\
    \bottomrule
  \end{tabular}
\end{table}

\subsection{AutoPhaseNN Reproduction}
The AutoPhaseNN reproduction is the anchor experiment. The expected narrative is:
the reproduced model should learn a compact object-domain amplitude/support
pattern and a phase field that remains spatially tied to the object region. The
paper should report both the far-field loss and the object-domain metrics, then
use the visualization in Fig.~\ref{fig:fig_autophasenn_reproduction.png} as the
qualitative baseline.

\todo{Insert final AutoPhaseNN checkpoint path, training epoch count, FT L1,
amplitude metrics, phase metrics, support metrics, and figure source.}

\subsection{UMamba Replacement}
The UMamba replacement is evaluated under the same data and physical
post-processing pipeline. The key observation to document is whether the
far-field metric improves without a corresponding improvement in amplitude,
phase, and support. If so, the result should be framed as an under-constrained
phase retrieval solution rather than simply as a failed optimizer run.

\todo{Insert UMamba checkpoint path, training epoch count, FT metrics, real-space
metrics, and representative visualization.}

\subsection{Diagnostic Ablation}
Only one ablation should be kept in the main paper: the one that most clearly
supports the explanation of the UMamba failure mode. For example, if the
soft-threshold experiment obtains acceptable numerical metrics but visibly
incorrect amplitude and phase, it can be used to show that smoother support
gradients alone do not solve the missing-prior problem. If a center-support
constraint is more informative, use that instead.

\todo{Choose one diagnostic ablation and report the same metric groups as the
other rows. Do not include every attempted training variant.}

\section{Discussion}

The experiments indicate that diffraction-domain loss alone can be insufficient
for judging reconstruction quality. A model may reduce FT L1 while producing an
incorrect object support or phase distribution. This is consistent with the
ill-posed nature of phase retrieval: multiple real-space configurations may
produce similar diffraction-domain errors unless additional priors are imposed.

\subsection{Why Far-field Fitting Can Be Misleading}
The far-field objective observes only the magnitude of the Fourier transform.
Consequently, it does not uniquely supervise the object-domain phase, support,
or spatial localization. The post-processing layer couples these variables, but
the loss is still applied after the FFT. A model can therefore exploit degrees
of freedom in amplitude and phase while still reducing the far-field residual.

\subsection{What the AutoPhaseNN Prior Provides}
The original AutoPhaseNN decoder appears to encode stronger spatial structure
through its architecture and post-processing assumptions. This prior may not be
a single explicit loss term. Instead, it can arise from the combination of
decoder topology, upsampling behavior, activation ranges, hard support masking,
and the oversampled diffraction setting. The practical effect is that the
network is biased toward compact object-domain solutions.

\subsection{Implication for Model Design}
The comparison suggests that replacing the backbone is not sufficient unless
the replacement preserves the physical bias of the original pipeline. A useful
future model should either inherit a hard spatial prior from its architecture or
enforce the prior through a simple post-processing constraint. Soft auxiliary
losses may help optimization, but they should not be presented as the central
solution unless they demonstrably improve real-space amplitude, phase, and
support at the same time.

\todo{Add concise analysis of the square/central support prior and oversampling
assumption, but keep the discussion focused on what is needed for the proposed
model design.}

\section{Conclusion}

This work frames the current experiments as a controlled study of neural phase
retrieval rather than only as a backbone replacement. Reproducing AutoPhaseNN
provides a reference pipeline, while the UMamba replacement tests whether a
different architecture can preserve the same physical behavior under aligned
training and evaluation conditions. The main open issue is the mismatch between
Fourier-domain fitting quality and real-space physical correctness. Future work
should focus on isolating the role of spatial support priors and designing hard
or architecture-level constraints that preserve physical plausibility without
overcomplicating the loss function.


\section*{Acknowledgment}
\todo{Add funding, lab, compute, and collaborator acknowledgments if needed.}

\bibliographystyle{IEEEtran}
\bibliography{references}

\end{document}
